\subsection{Two further congruence classes modulo five}

\begin{proposition}
\label{prop:mod-five-splitting}
Every positive exponent triple that is a permutation of $(1,1,2)$ or
$(1,4,4)$ modulo five gives a nonintegral ideal intersection graph.
\end{proposition}

\begin{proof}
The monic integer quintic $h_C(x)=\det(xI-C)/x$ is symmetric in the
exponents. Reducing the six-support quotient gives the complete identities
\begin{align*}
 (1,1,2):\quad h_C(x)&\equiv(x-1)^2(x-3)(x^2+3x+4),\\
 (1,4,4):\quad h_C(x)&\equiv(x-1)^2(x^3+x^2+4x+3)
 \pmod5.
\end{align*}
The quadratic discriminant is $3$, a nonsquare modulo five. The cubic
takes the values $3,4,3,1,4$ at $0,1,2,3,4$, respectively, and hence
is irreducible over $\mathbb F_5$. An integer quotient spectrum would
make $h_C$ split into linear factors modulo every prime, contradicting
either reduction. The complement and universal-class lifts give the
graph conclusion. No size, parity or endpoint-zero hypothesis is needed.
\end{proof}

On the equality subfamily $c=b(b+2)-a$, this gives a splitting obstruction
that a square-discriminant congruence can miss. If $b\equiv1\pmod5$, put
$q=ac\equiv a(3-a)$. The curve equation in
Remark~\ref{rem:endpoint-one-equality-curve} becomes
$G\equiv2(q-1)(q-2)\pmod5$. Thus $G=0$ permits precisely
$a\equiv1,2,4\pmod5$, giving one of the two excluded exponent classes.
For $q\equiv1$, the irreducible cubic has discriminant $1$ modulo five.
Consequently an integral equality spectrum requires $b\not\equiv1\pmod5$,
even though the discriminant-square test alone permits one excluded row.

\subsection{A local discriminant obstruction on the equality curve}

\begin{proposition}
\label{prop:equality-local-discriminant}
Let $a\leq b\leq c$ be positive integers with $c=b(b+2)-a$ and
$G(a,b)=0$, where $G$ is the equality equation in
Remark~\ref{rem:endpoint-one-equality-curve}. If an odd prime $p$ divides
$b$, $a^2\equiv1\pmod p$, and $13$ is a quadratic nonresidue modulo $p$,
then the ideal intersection graph is Laplacian nonintegral.
In particular, an integral equality spectrum requires $7\nmid b$ and
$11\nmid b$.
\end{proposition}

\begin{proof}
Put $q=ac$ and let $\Delta(q,b)$ be the discriminant of the explicit
cubic $P_3$ in Section~\ref{sec:open}. On $G=0$,
$h_C=(x-1)(x-b)P_3$, so an integer quotient spectrum would make
$\Delta$ an integer square. For $v=-1+2b+3b^2$, expansion of the
displayed coefficient formulas gives identities in $\mathbb Z[b,q]$:
\begin{align*}
 G(v,b)&=b^3(-b^4-8b^3+8b^2+20b-7),\\
 \Delta(v,b)&=13b^2+b^3R(b),\qquad R\in\mathbb Z[b],\\
 G(q,b)-G(v,b)&=(q-v)\bigl[(3b-1)(q+v)+b(b^2+4b-1)\bigr].
\end{align*}
The assumption $a^2\equiv1$ gives $q\equiv-1\pmod p$. Write
$e=v_p(b)\geq1$. The bracket in the third identity is $2$ modulo $p$,
so it is a unit. Since $G(q,b)=0$, the first and third identities imply
$q\equiv v\pmod{p^{3e}}$. The discriminant is an integer polynomial,
and therefore
\[
 \Delta(q,b)\equiv13b^2\pmod{p^{3e}}.
\]
Because $13$ is a nonresidue, $p\ne13$. Hence $v_p(\Delta)=2e$,
and the unit $\Delta/p^{2e}$ is congruent to
$13(b/p^e)^2$ modulo $p$. This is a nonresidue, contradicting square
discriminant. The complement and universal-class lifts give the graph
conclusion.

For $p=7$ or $11$, reducing the curve at $p\mid b$ gives
$q^2\equiv1$ and hence $a^4\equiv1$. Since $-1$ is a nonresidue at
both primes, $a^2\equiv1$ follows. Moreover $13$ reduces to the
nonresidues $6$ and $2$, respectively, so the first assertion applies.
\end{proof}

The obstruction resolves a repeated-root reduction: on the branch
$q\equiv-1$, one has $P_3\equiv x(x+1)^2\pmod p$ and
$\Delta\equiv0\pmod p$. Normalizing by $p^{2v_p(b)}$ reveals the
nonsquare unit that the reduction modulo $p$ misses. At $p=5$, this
excludes $(a,b)\equiv(1,0),(4,0)$ among the previously compatible
equality rows. The argument is uniform in $p$ and its valuation and
requires no finite exponent base.
